
En
effet, soient @@M0@@ l'espace affine @@M1@@ dont les @@M2@@-points correspondent aux suites @@M3@@, et @@M4@@ l'ouvert de @@M5@@ correspondant aux suites @@M6@@-régulières. Pour @@M7@@ fixé, les faisceaux @@M8@@ qui correspondent aux @@M9@@-points de @@M10@@ sont contenus dans une famille plate sur @@M11@@, et par hypothèse @@M12@@ est non-vide. Donc (i) résulte de la propriété de semi-continuité supérieure de la fonction @@M13@@ (EGA III, 7.7.5). Enfin, comme toute suite ``assez générale'' est @@M14@@-régulière et telle que @@M15@@ soit un sous-faisceau de @@M16@@ pour tout @@M17@@, (ii) résulte de (i).

\textbf{Lemme 1.7.} \emph{Soit @@M18@@ une suite exacte sur @@M19@@. Si le polynôme de Hilbert de @@M20@@ est écrit sous la forme}
@@M21@@
\emph{alors}
@@M22@@
En effet, @@M23@@; donc @@M24@@, d'où l'assertion.

\textbf{Proposition 1.8.} \emph{Soient @@M25@@ un @@M26@@-faisceau sur @@M27@@, @@M28@@ et @@M29@@ le polynôme de Hilbert de @@M30@@. Alors}~:
\begin{enumerate}[label=(\roman*)]
\item \emph{Pour @@M31@@, @@M32@@.}
\item \emph{@@M33@@ et @@M34@@ est aussi un @@M35@@-faisceau.}
\end{enumerate}
En effet, raisonnons par récurrence sur @@M36@@. Lorsque @@M37@@, c'est clair car @@M38@@. Lorsque @@M39@@, on a une suite exacte
@@M40@@
